\documentclass{article}
\usepackage{amsmath,amssymb}
\usepackage[margin=1in]{geometry}

\begin{document}

If
\[
f(x)
\]
is smooth over the mass of the Gaussian, then the integral
\[
I = \int_{\mathbb R^2} e^{-f(x)} \, \mathcal N(x \mid \mu, \Sigma)\, dx
\]
is essentially a Gaussian expectation:
\[
I = \mathbb E_{X \sim \mathcal N(\mu,\Sigma)}\!\left[e^{-f(X)}\right].
\]

The named methods that are commonly used fall into a few buckets:

\section*{(1) Laplace / saddlepoint approximations}
Write the exponent as
\[
g(x) = f(x) + \tfrac12 (x-\mu)^\top \Sigma^{-1}(x-\mu) + \text{const}.
\]
Then approximate $g$ by a quadratic around its minimizer $x^*$ (the mode of the unnormalized integrand). This gives a closed-form Gaussian integral and is extremely fast in 2D.

\paragraph{Keywords}
\begin{itemize}
  \item Laplace approximation
  \item saddlepoint approximation
  \item method of steepest descent
\end{itemize}

\paragraph{Best use}
When the integrand is unimodal and close to quadratic over the region that matters.

\section*{2) Delta method / Gaussian Taylor moment expansion}
Expand \(f(X)\) around \(\mu\):
\[
f(X) \approx f(\mu) + \nabla f(\mu)^\top (X-\mu) + \tfrac12 (X-\mu)^\top H(\mu)(X-\mu).
\]
Then \(e^{-f(X)}\) becomes the exponential of an affine-quadratic form in a Gaussian, which can be integrated in closed form if the resulting matrix is well-behaved.

\paragraph{Keywords}
\begin{itemize}
  \item delta method
  \item second-order delta method
  \item Gaussian moment expansion
  \item Taylor approximation under Gaussian
\end{itemize}

\section*{3) Gauss--Hermite quadrature in 2D}
Whiten the Gaussian by writing
\[
X = \mu + Lz, \qquad LL^\top = \Sigma,\qquad z \sim \mathcal N(0,I).
\]
Then
\[
I = \mathbb E_{z \sim \mathcal N(0,I)}\!\left[e^{-f(\mu + Lz)}\right].
\]
In 2D, tensor-product Gauss--Hermite quadrature is often very accurate and efficient.

\paragraph{Keywords}
\begin{itemize}
  \item Gauss--Hermite quadrature
  \item Hermite cubature
  \item Gaussian quadrature for normal expectations
\end{itemize}

\section*{4) Bayesian quadrature / Gaussian-process quadrature}
If function evaluations are expensive, Bayesian quadrature models the integrand with a Gaussian process and computes the Gaussian-weighted integral analytically under the posterior.

\paragraph{Keywords}
\begin{itemize}
  \item Bayesian quadrature
  \item Gaussian process quadrature
  \item probabilistic numerics integration
  \item WSABI (Warped Sequential Active Bayesian Integration)
\end{itemize}

\section*{5) Unscented transform / sigma-point methods}
These methods use a small deterministic set of points to approximate Gaussian expectations.

\paragraph{Keywords}
\begin{itemize}
  \item unscented transform
  \item sigma-point quadrature
  \item cubature Kalman filter
  \item spherical-radial cubature
\end{itemize}

\section*{6) Importance sampling / Monte Carlo}
Plain Monte Carlo is straightforward, but in 2D deterministic methods are usually preferable if \(f\) is smooth.

\section*{Rule of thumb}
\begin{itemize}
  \item \textbf{Fastest analytic approximation:} Laplace or second-order Taylor expansion.
  \item \textbf{High accuracy in 2D:} Gauss--Hermite quadrature.
  \item \textbf{Expensive function evaluations:} Bayesian quadrature / WSABI.
\end{itemize}

\section*{What to try first}
If \(f\) is smooth and 2D, a very practical first choice is:
\begin{itemize}
  \item whiten the Gaussian,
  \item apply 7--11 point Gauss--Hermite quadrature per dimension,
  \item compare against Laplace approximation around the mode.
\end{itemize}

\end{document}
